\section{Conclusion}
\label{sec:conclusion}

This paper set out two documents in one. The first is a description, at reimplementation depth, of
\Flux as it runs today: a consensus layer that gates block production on collateral rather than
hashpower, an orchestration layer that places 1,195 deployed applications across 6,489 collateralised
nodes by opportunistic self-selection (§3, \S\ref{sec:evaluation}), and the addressing, attestation and settlement
machinery that keeps that placement reachable and paid. The second is a specification, at the same
depth, of the architecture that follows if autonomous callers become the network's principal tenant:
demand-funded compensation, delegated agent authority, transparent settlement after the shielded
pools are retired, bounded and publicly verifiable bridging, and a hardened attestation substrate (Part IV). The two halves are written to the same standard: every
element of the second that is undecided is named as an open parameter or an open problem (\S\ref{sec:limitations}) rather
than filled with a plausible default.

The thesis under test, stated in §1, is a bet, not a finding: that autonomous software becomes the
dominant consumer of cloud capacity, and that \Flux's redesign follows from taking that bet seriously.
This paper does not argue the bet itself. What it argues is conditional and checkable --- given the
bet, whether the remaining work is extension of what already exists or invention of what does not ---
and \S\ref{sec:feasibility} answers it by classifying every mechanism in Part IV against that distinction: of twenty-two
assessed, six are Extension, eleven are Engineering, one moves from Research toward Engineering, and
four remain Research (\S\ref{sec:feasibility}, tally). Most of what the bet requires, this paper finds, \Flux can build by
extending components already in production. Four things it cannot.

The four are these, without re-arguing them here: post-quantum migration of outputs the protocol
cannot move; bounded agent spending on a chain whose script cannot observe history; unlinkable payment
for a specific deployment; and censorship resistance under measured collateral concentration
(\S\ref{sec:feasibility}, \S\ref{sec:limitations}).

One finding among these is the most consequential for how the remaining work should be ordered.
\PNR, as specified, would pay every node through the \PoN rotation independent of what that node
hosts, so the eligibility gate --- not the payment --- is what compels hosting (§7). That gate is
presently bound to neither the operator who ran the benchmark nor the machine that produced the
attestation (§4, §14). Activating \PNR before that binding exists would raise the payoff to passing
the gate without doing the work, by exactly what \PNR pays; hardening attestation is therefore a
precondition for \PNR, not a parallel workstream or a follow-on (\S\ref{sec:feasibility}, \S\ref{sec:migration}).

The argument is falsifiable in two independent places. The bet itself fails if, once the delegated
authority and machine-facing interfaces this paper specifies are wired and enforced (§17, §18), the
share of deployments attributable to autonomous callers rather than interactive operators stays
negligible --- the redesign would then be solving a problem the network does not have, and no property
argued in Parts II--IV would rescue that. The sequencing claim above fails on a narrower and
nearer-term observation: if \PNR activates while the benchmark and attestation signatures remain
unbound (§4, §14), and the network subsequently observes eligible nodes growing faster than any measure
of what they host, that is exactly the gaming this paper's ordering was written to prevent, and it
would show the priority argument of \S\ref{sec:migration}, not merely one component, to be wrong.

The open markers scattered through this document --- open parameters, open problems, and the notes
marked \emph{not established by this survey} --- are not a defect awaiting a final edit; they record decisions that a human must make and
that this paper deliberately left to a human rather than making silently. The next version of this
document should be judged by how many of them closed and whether closing them changed the mechanisms
it specifies, not by whether their count reached zero.

%% Drain this section's float queue before the next begins.
\FloatBarrier
